\documentclass[a4paper,12pt]{article}

\usepackage[T2A]{fontenc}
\usepackage[utf8]{inputenc}
\usepackage[russian]{babel}
\usepackage{amsmath,amssymb,mathtools}
\usepackage{PTSans}
\renewcommand{\familydefault}{\sfdefault}
\usepackage{icomma}
\usepackage[a4paper,margin=20mm,headheight=25pt]{geometry}
\usepackage{microtype}
\usepackage{tikz}
\usetikzlibrary{arrows.meta,positioning,shapes.geometric}
\usepackage{tabularx,booktabs,array}
\usepackage{enumitem}
\usepackage{needspace}
\usepackage{parskip}
\usepackage{fancyhdr}
\usepackage{setspace}
\usepackage{xcolor}

\definecolor{accent}{HTML}{008080}
\definecolor{darkaccent}{HTML}{004D4D}
\definecolor{textgray}{HTML}{333333}
\definecolor{softgray}{HTML}{6B7474}
\definecolor{rulegray}{HTML}{C7D1D1}

\color{textgray}
\onehalfspacing
\setlength{\parindent}{0pt}
\setlength{\parskip}{0.55em}
\setlength{\emergencystretch}{2em}
\widowpenalty=10000
\clubpenalty=10000
\displaywidowpenalty=10000

\setlist[itemize]{leftmargin=1.7em,itemsep=0.25em,topsep=0.35em}
\setlist[enumerate]{leftmargin=1.9em,itemsep=0.35em,topsep=0.35em}

\pagestyle{fancy}
\fancyhf{}
\fancyhead[L]{\small\textcolor{softgray}{Текстовые задачи на движение}}
\fancyhead[R]{\small\textcolor{softgray}{Анастасия Павлова}}
\fancyfoot[C]{\small\textcolor{softgray}{\thepage}}
\renewcommand{\headrulewidth}{0.4pt}
\renewcommand{\headrule}{%
  \hbox to\headwidth{\color{rulegray}\leaders\hrule height \headrulewidth\hfill}}

\newcommand{\sectiontitle}[1]{%
  \Needspace{6\baselineskip}%
  \vspace{0.25em}%
  {\Large\bfseries\color{darkaccent}#1}\par
  \vspace{0.25em}%
  \textcolor{rulegray}{\rule{\linewidth}{0.6pt}}%
  \vspace{0.35em}%
}

\newcommand{\subsectiontitle}[1]{%
  \Needspace{4\baselineskip}%
  \vspace{0.45em}%
  {\large\bfseries\color{darkaccent}#1}\par
  \vspace{0.15em}%
}

\newcommand{\important}[1]{%
  \par\vspace{0.2em}%
  \noindent
  \textcolor{accent}{\rule{2.4pt}{\baselineskip}}%
  \hspace{0.65em}%
  \begin{minipage}[t]{0.94\linewidth}
  #1
  \end{minipage}%
  \par\vspace{0.3em}%
}

\tikzset{
  flowstart/.style={
    ellipse,
    draw=darkaccent,
    line width=0.9pt,
    align=center,
    minimum width=38mm,
    minimum height=9mm,
    inner sep=2.5mm
  },
  flowaction/.style={
    rounded corners=3pt,
    rectangle,
    draw=accent,
    line width=0.9pt,
    align=center,
    text width=118mm,
    minimum height=11mm,
    inner sep=2.8mm,
    font=\small
  },
  flowarrow/.style={
    -{Stealth[length=3mm,width=2mm]},
    line width=0.9pt,
    draw=textgray
  }
}

\begin{document}

% ============================================================
% Титульный лист
% ============================================================
\begin{titlepage}
\thispagestyle{empty}
\centering

\vspace*{2.15cm}

{\large\textcolor{softgray}{Математика \(\cdot\) подготовка к ЕГЭ}\par}

\vspace{1.15cm}

{\Huge\bfseries\color{darkaccent}Чек-лист\par}

\vspace{0.45cm}

{\LARGE\bfseries Текстовые задачи на движение\par}

\vspace{0.6cm}

{\large
встречное движение \(\cdot\) неодновременный старт \(\cdot\)\par
движение в догонку \(\cdot\) разворот \(\cdot\) единицы измерения
\par}

\vspace{2.1cm}



\vfill

{\normalsize Лёвин Артём Александрович\par}
{\normalsize эксклюзивно для Анастасии Павловой\par}

\vspace{0.55cm}

{\normalsize \today\par}

\vspace{1.7cm}

\begin{tikzpicture}[remember picture,overlay]
  \draw[accent!48,line width=1.15pt]
    ([xshift=18mm,yshift=18mm]current page.south west)
    .. controls
    ([xshift=67mm,yshift=31mm]current page.south west)
    and
    ([xshift=-72mm,yshift=8mm]current page.south east)
    ..
    ([xshift=-18mm,yshift=22mm]current page.south east);
\end{tikzpicture}

\end{titlepage}

% ============================================================
% Основной чек-лист
% ============================================================
\sectiontitle{1. Основа: таблица \(v\)--\(t\)--\(S\)}

Для равномерного движения используются три связанные величины:
\[
\boxed{S=vt},
\qquad
\boxed{v=\frac{S}{t}},
\qquad
\boxed{t=\frac{S}{v}}.
\]

Здесь \(S\) --- пройденный путь, \(v\) --- скорость, \(t\) --- время движения.

По смыслу задач на движение обычно выполняются ограничения
\[
t>0,\qquad v>0,\qquad S\ge 0.
\]
Если \(x\) обозначает часть ограниченного отрезка пути, дополнительно учитывайте соответствующий диапазон значений.

Удобно располагать данные в порядке
\[
\boxed{v\qquad t\qquad S},
\]
поскольку первые две величины непосредственно дают третью:
\[
\boxed{v\cdot t=S}.
\]

\begin{center}
\renewcommand{\arraystretch}{1.35}
\begin{tabular}{@{}ccc@{}}
\toprule
Скорость & Время & Путь\\
\midrule
\(v\) & \(t\) & \(S\)\\
\bottomrule
\end{tabular}
\end{center}

\important{
\textbf{Главный принцип заполнения таблицы.}
Сначала обозначьте через \(x\) то, что требуется найти, затем внесите данные из условия.
Оставшийся столбец заполните по формуле движения.
}

\subsectiontitle{Шесть шагов решения}

\begin{enumerate}
  \item Определите, что требуется найти.
  \item При необходимости сделайте короткую схему движения.
  \item Введите переменную \(x\) и выразите через неё связанные величины.
  \item Заполните известными данными таблицу \(v\)--\(t\)--\(S\).
  \item Недостающий столбец получите из формулы \(S=vt\), \(t=S/v\) или \(v=S/t\).
  \item Используйте оставшееся смысловое условие, составьте уравнение, решите его и проверьте единицы измерения.
\end{enumerate}

\important{
Уравнение должно выражать реальную связь из условия: сумму путей, равенство времён,
разницу времён, начальный отрыв или другую конкретную величину.
}

\subsectiontitle{Как выбирать \(x\)}

Если требуется найти время до встречи, удобно положить
\[
t=x.
\]
Если требуется скорость первого автомобиля,
\[
v_1=x.
\]
Если требуется путь первого автомобиля,
\[
S_1=x.
\]

Иногда искомая величина является только частью маршрута. Тогда весь путь участника
нужно выразить через \(x\). Например, при развороте у опушки пути могут иметь вид
\[
6+x
\qquad\text{и}\qquad
6-x.
\]

% ============================================================
% Встречное движение
% ============================================================
\sectiontitle{2. Встречное движение: одновременный старт}

Два автомобиля одновременно выехали навстречу друг другу из городов,
расстояние между которыми равно \(560\) км. Скорости равны \(60\) км/ч и \(80\) км/ч.
Найдём время до встречи.

Поскольку автомобили стартовали одновременно и рассматривается один и тот же момент встречи,
\[
t_1=t_2=x.
\]

\begin{center}
\renewcommand{\arraystretch}{1.35}
\begin{tabular}{@{}lccc@{}}
\toprule
 & \(v\), км/ч & \(t\), ч & \(S\), км\\
\midrule
Первый автомобиль & \(60\) & \(x\) & \(60x\)\\
Второй автомобиль & \(80\) & \(x\) & \(80x\)\\
\bottomrule
\end{tabular}
\end{center}

К моменту встречи сумма пройденных путей равна расстоянию между городами:
\[
60x+80x=560.
\]
Отсюда
\[
140x=560,
\qquad
\boxed{x=4\text{ ч}}.
\]

\important{
При движении навстречу друг другу до первой встречи обычно выполняется
\[
\boxed{S_1+S_2=S_{\text{общее}}}.
\]
}

\subsectiontitle{Если требуется найти место встречи}

Пусть расстояние между городами равно \(480\) км, а скорости автомобилей ---
\(55\) км/ч и \(65\) км/ч. Требуется найти путь первого автомобиля до встречи.

Положим
\[
S_1=x,
\qquad
S_2=480-x,
\qquad
0<x<480.
\]
Так как автомобили выехали одновременно,
\[
\frac{x}{55}=\frac{480-x}{65}.
\]
Сократим вычисления до перемножения крупных чисел:
\[
13x=11(480-x),
\]
\[
24x=5280,
\qquad
\boxed{x=220\text{ км}}.
\]

\subsectiontitle{Если требуется найти скорость}

Из двух городов, расстояние между которыми равно \(390\) км, одновременно навстречу
друг другу выехали два автомобиля и встретились через \(3\) часа.
Скорость второго автомобиля равна \(60\) км/ч. Найдём скорость первого.

Пусть
\[
v_1=x.
\]
Тогда за \(3\) часа автомобили прошли соответственно
\[
3x
\qquad\text{и}\qquad
3\cdot60=180
\]
километров. Поэтому
\[
3x+180=390,
\]
\[
3x=210,
\qquad
\boxed{x=70\text{ км/ч}}.
\]

% ============================================================
% Неодновременный старт
% ============================================================
\sectiontitle{3. Встречное движение: старт в разное время}

Из города \(A\) в город \(B\) со скоростью \(80\) км/ч выехал первый автомобиль.
Через \(2\) часа навстречу ему из города \(B\) выехал второй автомобиль
со скоростью \(60\) км/ч. Расстояние между городами равно \(580\) км.
Найдём, через сколько часов после выезда второго автомобиля произойдёт встреча.

Пусть второй автомобиль двигался до встречи \(x\) часов.
Первый выехал на \(2\) часа раньше, поэтому его время движения равно
\[
x+2.
\]

\begin{center}
\renewcommand{\arraystretch}{1.35}
\begin{tabular}{@{}lccc@{}}
\toprule
 & \(v\), км/ч & \(t\), ч & \(S\), км\\
\midrule
Первый автомобиль & \(80\) & \(x+2\) & \(80(x+2)\)\\
Второй автомобиль & \(60\) & \(x\) & \(60x\)\\
\bottomrule
\end{tabular}
\end{center}

Сумма путей к моменту встречи равна \(580\) км:
\[
80(x+2)+60x=580.
\]
\[
140x+160=580,
\qquad
140x=420,
\qquad
\boxed{x=3\text{ ч}}.
\]

\important{
\textbf{Ключевая мысль:} кто выехал раньше, тот к одному и тому же моменту двигался дольше.
Если первый участник стартовал на \(a\) часов раньше второго, удобно записывать
\[
t_2=x,
\qquad
t_1=x+a.
\]
}

\subsectiontitle{Когда разность времён входит прямо в уравнение}

Если пути и скорости уже выражены, время можно найти по формуле
\[
t=\frac{S}{v}.
\]
Если первый участник выехал на \(1\) час раньше второго, то
\[
\boxed{t_1-t_2=1}.
\]
Например, если при встречном движении общее расстояние равно \(380\) км,
пути участников равны \(x\) и \(380-x\) км, а скорости --- \(50\) и \(60\) км/ч, то
\[
\frac{x}{50}-\frac{380-x}{60}=1.
\]
Знак разности выбирается по смыслу: из большего времени вычитается меньшее.

% ============================================================
% Блок-схема
% ============================================================
\clearpage
\thispagestyle{empty}

\begin{center}
{\Large\bfseries\color{darkaccent}Алгоритм решения задачи на движение}
\end{center}

\vspace{8mm}

\begin{center}
\begin{tikzpicture}[node distance=6.5mm]

\node[flowstart] (start) {Начало};

\node[flowaction,below=of start] (step1)
{1. Прочитайте условие и определите, какую величину требуется найти.};

\node[flowaction,below=of step1] (step2)
{2. Разберите схему движения: направления, моменты старта, встречу, догонку или разворот. При необходимости сделайте чертёж.};

\node[flowaction,below=of step2] (step3)
{3. Введите \(x\) для искомой величины и выразите через \(x\) связанные с ней данные.};

\node[flowaction,below=of step3] (step4)
{4. Заполните таблицу \(v\)--\(t\)--\(S\): сначала известные данные, затем недостающий столбец по формуле движения.};

\node[flowaction,below=of step4] (step5)
{5. Переведите оставшееся смысловое условие в уравнение: сумма путей, равенство или разность времён, начальный отрыв.};

\node[flowaction,below=of step5] (step6)
{6. Решите уравнение, проверьте единицы измерения, смысл результата и соответствие вопросу задачи.};

\node[flowstart,below=of step6] (finish) {Ответ};

\draw[flowarrow] (start) -- (step1);
\draw[flowarrow] (step1) -- (step2);
\draw[flowarrow] (step2) -- (step3);
\draw[flowarrow] (step3) -- (step4);
\draw[flowarrow] (step4) -- (step5);
\draw[flowarrow] (step5) -- (step6);
\draw[flowarrow] (step6) -- (finish);

\end{tikzpicture}
\end{center}

\clearpage

% ============================================================
% Догонка
% ============================================================
\sectiontitle{4. Движение в догонку}

Пусть города \(A\), \(B\) и \(C\) расположены на одной дороге, причём \(B\) находится
между \(A\) и \(C\), а расстояние \(AB\) равно \(125\) км.
Из \(A\) в сторону \(C\) выехал легковой автомобиль, а одновременно из \(B\)
в ту же сторону выехал грузовик. Скорость легкового автомобиля на \(25\) км/ч больше.
Найдём время до догоняния.

Пусть скорость грузовика равна \(y\) км/ч. Тогда скорость легкового автомобиля равна
\[
y+25.
\]
Пусть до догоняния прошло \(x\) часов.

\begin{center}
\renewcommand{\arraystretch}{1.35}
\begin{tabular}{@{}lccc@{}}
\toprule
 & \(v\), км/ч & \(t\), ч & \(S\), км\\
\midrule
Легковой автомобиль & \(y+25\) & \(x\) & \(x(y+25)\)\\
Грузовик & \(y\) & \(x\) & \(xy\)\\
\bottomrule
\end{tabular}
\end{center}

К моменту догоняния легковой автомобиль прошёл на \(125\) км больше:
\[
x(y+25)=xy+125.
\]
После раскрытия скобок
\[
xy+25x=xy+125,
\]
поэтому
\[
25x=125,
\qquad
\boxed{x=5\text{ ч}}.
\]

Переменная \(y\) сократилась: время догоняния определяется начальным отрывом
и разностью скоростей, равной \(25\) км/ч.

\important{
Для движения в догонку удобно проверять смысл уравнения по формуле
\[
\boxed{S_{\text{быстрого}}-S_{\text{медленного}}
=S_{\text{начальный отрыв}}}.
\]
}

% ============================================================
% Единицы измерения
% ============================================================
\sectiontitle{5. Единицы измерения}

Перед составлением уравнения единицы должны быть согласованы.
Если скорость дана в км/ч, удобно выражать расстояние в километрах, а время --- в часах.

Например,
\[
300\text{ м}=0{,}3\text{ км}.
\]

Пусть быстрый пешеход идёт на \(2\) км/ч быстрее медленного, а расстояние между ними
составляет \(300\) м. Если оба идут в одном направлении, то за время \(x\) часов
быстрый должен пройти на \(0{,}3\) км больше:
\[
2x=0{,}3.
\]
Отсюда
\[
x=0{,}15\text{ ч}.
\]
Переведём в минуты:
\[
0{,}15\cdot60=9.
\]
Следовательно,
\[
\boxed{9\text{ мин}}.
\]

\begin{center}
\renewcommand{\arraystretch}{1.35}
\begin{tabularx}{0.88\linewidth}{@{}lX@{}}
\toprule
Переход & Формула\\
\midrule
метры \(\to\) километры & \(\displaystyle S_{\text{км}}=\frac{S_{\text{м}}}{1000}\)\\[2mm]
километры \(\to\) метры & \(\displaystyle S_{\text{м}}=1000S_{\text{км}}\)\\[2mm]
часы \(\to\) минуты & \(\displaystyle t_{\text{мин}}=60t_{\text{ч}}\)\\[2mm]
минуты \(\to\) часы & \(\displaystyle t_{\text{ч}}=\frac{t_{\text{мин}}}{60}\)\\
\bottomrule
\end{tabularx}
\end{center}

\important{
После решения обязательно проверьте, в каких единицах получен \(x\).
Если время находилось при скорости в км/ч и расстоянии в километрах, \(x\) первоначально выражен в часах.
}

% ============================================================
% Разворот
% ============================================================
\sectiontitle{6. Движение с разворотом}

Два пешехода одновременно вышли из одной точки к опушке,
расположенной в \(6\) км от старта. Их скорости равны \(6{,}5\) км/ч и \(3{,}5\) км/ч.
Быстрый пешеход дошёл до опушки, сразу развернулся и встретил второго.
Пусть место встречи находится на расстоянии \(x\) км от опушки. По смыслу
\[
0<x<6.
\]

Быстрый прошёл
\[
6+x
\]
километров, а медленный ---
\[
6-x
\]
километров.

Оба вышли одновременно и двигались до одного и того же момента встречи, поэтому их времена равны:
\[
\frac{6+x}{6{,}5}=\frac{6-x}{3{,}5}.
\]
Перемножим крест-накрест:
\[
3{,}5(6+x)=6{,}5(6-x).
\]
\[
21+3{,}5x=39-6{,}5x,
\]
\[
10x=18,
\qquad
x=1{,}8.
\]
Следовательно,
\[
\boxed{x=1{,}8\text{ км}=1800\text{ м}}.
\]

\important{
При развороте в таблице указывается весь путь, фактически пройденный участником.
Искомый \(x\) может быть только фрагментом этого пути.
}

% ============================================================
% Самопроверка
% ============================================================
\sectiontitle{7. Самопроверка и типичные ошибки}

Перед записью ответа проверьте:

\begin{enumerate}[itemsep=0.2em]
  \item Соответствует ли \(x\) вопросу задачи и все ли данные перенесены в таблицу?
  \item Учтены ли направления движения и моменты старта?
  \item Верно ли записана смысловая связь между участниками?
  \item При развороте указан ли весь фактически пройденный путь?
  \item Согласованы ли единицы измерения?
  \item Имеет ли результат смысл и отвечает ли он на вопрос задачи?
\end{enumerate}

\subsectiontitle{Типичные ошибки}

\begin{itemize}[itemsep=0.15em]
  \item Одинаковое время у участников, стартовавших в разные моменты.
  \item \(x\) принят за весь путь, хотя обозначает только его часть.
  \item При догонке вместо разности путей использована их сумма.
  \item Смешаны метры с километрами или минуты с часами.
\end{itemize}

\important{
Короткая схема рассуждения:
\[
\boxed{\text{текст}\to\text{схема}\to\text{таблица}\to\text{формула}\to\text{уравнение}\to\text{проверка}}.
\]
}

% ============================================================
% Тренировочный блок
% ============================================================
\clearpage
\sectiontitle{Тренировочный блок. Путь А}

\textbf{5 заданий.} Решайте с помощью таблицы \(v\)--\(t\)--\(S\).
Задачи расположены по нарастающей сложности.

\begin{enumerate}[label=\textbf{\arabic*.},leftmargin=2.1em,itemsep=0.85em,topsep=0.4em]
  \item Из двух городов, расстояние между которыми равно \(640\) км,
  одновременно навстречу друг другу выехали два автомобиля со скоростями
  \(70\) км/ч и \(90\) км/ч. Через сколько часов они встретятся?

  \item Из двух городов, расстояние между которыми равно \(540\) км,
  одновременно навстречу друг другу выехали два автомобиля со скоростями
  \(60\) км/ч и \(75\) км/ч. На каком расстоянии от города, из которого выехал
  первый автомобиль, произойдёт встреча?

  \item Из города \(A\) в город \(B\) со скоростью \(80\) км/ч выехал автомобиль.
  Через \(2\) часа навстречу ему из города \(B\) выехал второй автомобиль
  со скоростью \(60\) км/ч. Расстояние между городами равно \(580\) км.
  Через сколько часов после выезда второго автомобиля произойдёт встреча?

  \item По одной дороге в одном направлении одновременно начали движение два пешехода.
  Быстрый находится позади медленного на \(450\) м. Скорость медленного равна
  \(5\) км/ч, а скорость быстрого на \(3\) км/ч больше.
  Через сколько минут быстрый пешеход догонит медленного?

  \item Два пешехода одновременно вышли из одной точки к опушке леса,
  находящейся в \(6\) км от старта. Их скорости равны \(6{,}5\) км/ч и \(3{,}5\) км/ч.
  Первый дошёл до опушки, сразу развернулся и встретил второго.
  На каком расстоянии от опушки произошла встреча? Ответ дайте в метрах.
\end{enumerate}

% ============================================================
% Ответы
% ============================================================
\clearpage
\sectiontitle{Ответы}

\begin{center}
\renewcommand{\arraystretch}{1.45}
\begin{tabularx}{0.72\linewidth}{@{}>{\bfseries}cX@{}}
\toprule
№ & Ответ\\
\midrule
1 & \(4\) ч\\
2 & \(240\) км\\
3 & \(3\) ч\\
4 & \(9\) мин\\
5 & \(1800\) м\\
\bottomrule
\end{tabularx}
\end{center}

\vfill

\begin{center}
\textcolor{softgray}{\small
Лёвин Артём Александрович \(\cdot\) эксклюзивно для Анастасии Павловой}
\end{center}

\end{document}